\section{Case-Study 3: Fish Stunning - Erfolgsquotenbestimmung} \label{sec:case_study_3}
Im Rahmen der vorliegenden Arbeit wird eine Fallstudie herangezogen, 
um die entwickelte Strategie zur Risikobehandlung in einem praxisnahen Anwendungskontext zu demonstrieren. 
Die Fallstudie basiert auf einem Projekt aus dem industriellen Umfeld der 
Firma Baader sowie deren Tochtergesellschaft TriVision. 
Aufgrund eingeschränkter Verfügbarkeit projektspezifischer Dokumentation 
und fehlender direkter Projektbeteiligung des Autors erhebt die Darstellung 
keinen Anspruch auf vollständige empirische Abbildung. 
Stattdessen wird der Anwendungsfall in abstrahierter und teilweise 
hypothetisch ergänzter Form rekonstruiert, um eine konsistente und für die Analyse geeignete Ausgangsbasis zu schaffen.

Gegenstand des betrachteten Projekts ist die Entwicklung eines 
optischen Erkennungssystems zur Bewertung des Betäubungserfolgs von Fischen 
in einer industriellen Verarbeitungsanlage. 
Konkret soll auf Basis von Live-Bilddaten einer Fisch Verarbeitungsmaschine automatisiert bestimmt werden, 
ob ein Fisch ordnungsgemäß ,,gestunnt'' wurde, bevor der Prozess fortgeführt wird. Zielgröße des Systems ist die Ermittlung einer Erfolgsquote, 
welche insbesondere vor dem Hintergrund tierschutzrechtlicher Anforderungen von zentraler Bedeutung ist. 
Wird eine definierte Mindestquote unterschritten, sind entsprechende Gegenmaßnahmen einzuleiten, 
beispielsweise in Form technischer Wartung oder Anpassung der Prozessparameter.

Die technische Umsetzung erfolgt mittels eines neuronalen Netzes aus 
dem Bereich des Maschinelles Lernen, das unter Verwendung der Softwarebibliothek TensorFlow entwickelt wurde. 
Für den produktiven Einsatz wird das Modell in das interoperable Open Neural Network Exchange-Format überführt, 
um eine Ausführung auf spezialisierter Zielhardware zu ermöglichen. 
Der Entwicklungsprozess kann dabei nicht als streng methodisch oder agil strukturiert beschrieben werden, 
sondern folgte eher einem iterativen „Trial-and-Error“-Ansatz, was zusätzliche Unsicherheiten 
hinsichtlich Modellqualität und Systemstabilität impliziert. Eine abschließende wiederkehrende 
Modellevaluation ist nicht vorgesehen.

Vor dem Hintergrund dieses Anwendungsszenarios ergeben sich verschiedene potenzielle Risikofelder. 
Hierzu zählen insbesondere datengetriebene Risiken wie ein möglicher 
Data Drift infolge veränderter Umweltbedingungen (z. B. Lichtverhältnisse), der die Modellperformance 
negativ beeinflussen kann. Darüber hinaus bestehen erhebliche ethische Risiken, 
insbesondere im Fall von Fehlklassifikationen, bei denen nicht korrekt betäubte Tiere fälschlicherweise 
als betäubt erkannt werden. Ergänzend sind auch betriebswirtschaftliche und reputationsbezogene Risiken 
zu berücksichtigen, etwa wenn das System beim Kunden nicht zuverlässig operationalisiert 
werden kann oder die erwartete Leistungsfähigkeit nicht erreicht wird.

Die nachfolgende Analyse nutzt dieses Szenario als strukturierte Grundlage, 
um die entwickelte Strategie zur systematischen Identifikation, 
Bewertung und Behandlung von Risiken exemplarisch anzuwenden und deren praktische Umsetzbarkeit zu demonstrieren.

\subsection{Semi-quantitative Bewertungsstruktur}

Für die Durchführung des Gedankenexperiments wird eine semi-quantitative Bewertungsstruktur 
zur Risikobewertung verwendet. Die Bewertung erfolgt auf einer ordinalen 
Skala von 1 bis 5 und orientiert sich konzeptionell an der \textit{Fehlermöglichkeits- und -einflussanalyse} (FMEA), 
wird jedoch gezielt erweitert.

Im Gegensatz zur klassischen FMEA, die typischerweise auf einer multiplikativen 
Verknüpfung der Bewertungsdimensionen basiert, 
wird in der vorliegenden Arbeit ein gewichtetes additives Modell eingesetzt. 
Ziel dieser Anpassung ist es, eine unverhältnismäßige Skalierung der Risikowerte 
zu vermeiden und gleichzeitig eine transparente Priorisierung der einzelnen Bewertungsdimensionen zu ermöglichen.

Die Bewertungsstruktur umfasst die folgenden drei Kennziffern:
\begin{itemize}
    \item Eintrittswahrscheinlichkeit (1 = unwahrscheinlich, 5 = sehr wahrscheinlich)
    \item Auswirkung (1 = gering , 5 = stark)
    \item Entdeckbarkeit (1 = schwer, 5 = leicht)
\end{itemize}

Diese werden unterschiedlich gewichtet, um ihrer relativen Bedeutung im Anwendungskontext Rechnung zu tragen. 
Die Gewichtung ist in Tabelle \ref{tab:gewichtung} dargestellt.

\begin{table}[h]
\centering
\label{tab:gewichtung}
    \begin{tabular}{l c}
    \hline
    \textbf{Kennziffer} & \textbf{Gewichtung} \\
    \hline
    Eintrittswahrscheinlichkeit & 0{,}4 \\
    Auswirkung                  & 0{,}4 \\
    Entdeckbarkeit              & 0{,}2 \\
    \hline
    \end{tabular}
\caption{Gewichtung der Bewertungsdimensionen}
\end{table}

Die Aggregation der Einzelbewertungen zu einem Key Risk Indicator (KRI) erfolgt über folgende Formel:

\begin{equation}
KRI = 0{,}4 \cdot Eintrittswahrscheinlichkeit + 0{,}4 \cdot Auswirkung + 0{,}2 \cdot Entdeckbarkeit
\end{equation}

Durch diese additive Verknüpfung ergibt sich ein Wertebereich von minimal 1 bis maximal 5. 
Diese KRI's bilden die Entscheidungsgrundlage der Risk-Gates welche sich, 
wie in Tabelle \ref{tab:risikoklassen} darstellt.

\begin{table}[h]
\centering
\label{tab:risikoklassen}
    \begin{tabular}{c l l}
    \hline
    \textbf{KRI} & \textbf{Risikokategorie} & \textbf{Ampelstatus} \\
    \hline
    $1 \leq x < 2$ & geringes Risiko        & grün \\
    $2 \leq x < 4$ & erhöhtes Risiko       & gelb \\
    $4 \leq x \leq 5$ & hohes Risiko          & rot \\
    \hline
    \end{tabular}
\caption{Risikoklassen und Bewertungsskala}
\end{table}

Die gewählte Kategorisierung ist bewusst konservativ ausgestaltet. 
Insbesondere erfolgt die Einordnung in kritische Bereiche (gelb bzw. rot) vergleichsweise früh. 
Diese Entscheidung basiert auf der Annahme, dass im betrachteten Anwendungskontext 
ein hoher Fokus auf der schnellen Operationalisierung technischer Systeme liegt, 
wodurch Risiken potenziell unterbewertet werden könnten.

Die vorliegende Einteilung stellt somit einen bewussten methodischen Gegenpol dar, 
der eine frühzeitige Sensibilisierung für kritische Entwicklungen unterstützt und eine strengere 
Risikosteuerung im weiteren Verlauf des Entwicklungs- und Einführungsprozesses fördert.

Die Integration von Ergebnissen aus weiterführenden Analyseverfahren 
(z.\,B. Sensitivitätsanalysen) in die vorliegende Bewertungsstruktur erfolgt 
über geeignete Transformations- bzw. Mapping-Ansätze. 
Dabei werden methodenspezifische Kennzahlen in die einheitliche Risikoskala überführt. 

Für die Nachverfolgung der Risiken und späteren Gegenüberstellung, werden die KRI's samt 
ihrer Maßnahmen in einem Risk-Register verwaltet und zusammengeführt. Dieses bildet 
zusätzlich noch den KRI vorher und nachher ab, wodurch eine Verändung direkt sichtbar wird. 

Im folgenden Abschnitt werden die Phasen des DASC-PM gedanklich durchgeführt,
die möglichen Risiken identifiziert und in die eingangs vorgestellte Risikobewertung
überführt. Abschließend zu jeder Phase wird auf die entsprechende Entscheidung in 
dem korrespondierenden Risk-Gate eingegangen.

\subsection{Phase 1: Problem Statement}

    Mögliche Risiken dieser Phase: 
    - Keine besonders relevanten bzw. interessanten Risiken hier. Es werden 
    keine Personenbezogenen Daten erhoben, Das Ziel ist klar definiert, 
    Ressourcen werden von TriVision gestellt bzw. dort wird es umgesetzt,
    diese Anpassung ist gefordert und somit ist der Businesscase auch gegeben,

    \paragraph{Das Risk-Gate bleibt auf Grün} da hier keine relevanten Risiken identifiziert werden konnten.

\subsection{Phase 2: Data Provisioning}
    Mögliche Risiken dieser Phase: 
    - Ausgewähltes Kamerasystem zu schlecht und somit ist die Datenqualität nicht ausreichend. (Maßnahme -> Bessere Kamera)
        - Die erste Datensichtung hat gezeigt dass die Datenqualität der Videoaufnahmen zu schlecht sein könnte, um daraus 
          tatsächlich ein Modell zu trainieren.
        - (E)intrittswahrscheinlichkeit = 2 - Es könnte auch ausreichen und das Kamerasystem ist günstiger als das nächste Upgrade.
        - (A)uswirkung = 5 - Produkt funktioniert nicht zuverlässig
        - (Ent)deckbarkeit = 2 - Kein aktives Kontrollsystem, aber im weiteren Prozess würde es auffallen.
        KRI = 2 * 0.4 + 5 * 0.4 + 2 * 0.2 = 3.2 
    - Infrastruktur nicht ausreichend für die Datenmengen (Maßnahmen -> Videos nur kurz speichern und dann löschen, Direkte Verbindung von Modell zu Kamera)
        - Die entstandenen Aufnahmen sind sehr groß und die ersten Architekturkonzepte sahen einen Transport über das Netzwerk vor. Zudem sollten
        die Videos für eine definierte Zeit zur Nachverfolgung gespeichert werden.
        - E = 2 - Die erste Analyse hat gezeigt, dass es grundsätzlich möglich ist, aber es bestehen Sorgen, dass die Zukünftige Last das System überfordern könnte.
        - A = 5 - Das System würde zusammenbrechen und somit gar nicht funktionieren.
        - Ent = 5 - Fällt während der Entwicklung und ggf. in der Produktion auf.
        KRI = 2 * 0.4 + 5 * 0.4 + 5 * 0.2 = 3.8
    \paragraph{Dieser Abschnitt schließt mit einer gelben Ampel ab} und erlaubt ein Fortfahren in die nächste Phase. Es werden 
    lediglich zwei Risiken aufgenommen, die bis auf weiteres beobachtet werden müssen. 
    
\subsection{Phase 3: Analysis}
    Mögliche Risiken dieser Phase: 
    - Zuverlässigkeit der Ergebnisse (Analyse und KRI durch Performance Indicator und Modellkennziffern.)
        - Es besteht ein Ethisches Risiko, wenn die Ergebnisse nicht korrekt klassifizieren. Dadurch würden 
        Fische automatisiert als ,,betäubt'' gezählt werden und die Erfolgsquote würde fälschlicherweise 
        als hoch eingestuft werden. Dadurch würde Tierleid in Kauf genommen werden.
        Das Modell wurde anhand der bestehenden Videos trainiert, da keine neue Kamera besorgt wurde. 
        Die Trainingsdaten wurden gelabelt und waren in ausreichendem Umfang vorhanden.
        Die Kontrolle der Modellperformance wurde anhand der typischen Metriken durchgeführt.
        Im Anhang ist eine Liste mit den erhobenen Performance Parametern die durch eine 
        Schwellwertbasierte Einordnung in KRI's überführt wurden. Im Hauptmarkt für diese Maschine 
        wird eine Betäubung von 100\% angestrebt, es gibt aber keine absolute regulatorische Vorgabe dazu \citep{espmark_assessment_2025}. Dies wird regelmäßig kontrolliert und durch eine 
        Institution nachverfolgt. Fabriken müssen schließen, wenn dies nicht gewährleistet werden kann.
        - Folgende Metriken wurden weiter betrachtet und in Risiken überführt:
            - (FNR) Fehlklassifikation: Nicht betäubter Fisch wird als betäubt erkannt (kritischer False Negative Fall) 
                - Aus der ,,schlechten'' False Negative Quote von 25\% lässt sich eine hohe Eintrittswahrscheinlichkeit anleiten
                  E = 4
                  Der Kontext ist sehr schwerwiegend, da hier auch Tierwohl Bedenken bzw. regulatorische Quoten nicht mehr eingehalten werden können
                  A = 5
                  Diese Metriken werden bei jedem Training erzeugt und somit gibt es hier eine wahrscheinliche Entdeckung dieses Risikos.
                  Ent = 2
                  KRI = 4 * 0.4 + 5 * 0.4 + 2 * 0.2 = 4.0
            - (Recall) Unzureichende Erkennungsrate tatsächlich betäubter Fische (Recall-Schwäche)
                  Der Recall ist ebenfalls unzureichend da nur ca. 72\% der Fische korrekt positiv klassifiziert werden.
                  E = 4
                  Die Einordnung dieser Bewertung ist ähnlich zu dem vorherigen.
                  A = 5
                  Ent = 2
                  KRI = 4 * 0.4 + 5 * 0.4 + 2 * 0.2 = 4.0
            - (ROC-AUC) Geringe Trennschärfe des Modells (unsichere Klassifikation)
                  Der AUC Wert von 0.80 zeigt eine über alle Schwellwerte genaue gewichtung 
                  des ,,korrekten'' Falls bzw. Outcomes über den falschen. Das ist ebenfalls noch nicht gut genug.
                  E = 4
                  A = 5
                  Ent = 2
                  KRI = 4 * 0.4 + 5 * 0.4 + 2 * 0.2 = 4.0
    \paragraph{Das Risk-Gate zeigt nach dieser Phase eine rote Ampel} die sich aus den kritischen Performance Ergebnissen 
    der Entwicklung ergibt. Hier ergeben sich nun Maßnahmen, die die gesammelten Risiken addressieren.

    \notiz{TODO: Maßnahmen und Bearbeitung beschreiben -> Freigabe durch Verbesserung der Metriken -> Weiterführung des Prozesses}

\subsection{Phase 4: Deployment}
    Mögliche Risiken dieser Phase: 
    - Data-Drift durch z.B. sich ändernde Lichtverhältnisse (Jahreszeitenabhängiges Nebenlicht in Norwegen ist sehr unterschiedlich. Maßnahme: Hardware nachrüsten für Beleuchtung oder softwareseitig nachbelichten.)
        - (Performance Drop) Instabilität bei veränderten Umweltbedingungen (z.\,B. Lichtverhältnisse, Data Drift)
    - Andere Fischzüchtungen führen ggf. zu Modell-Drift da sie sich anders verhalten (Andere Zuchtbecken führen ggf. zu anderen Ergebnissen. Maßnahmen: Datenerhebung aus mehreren Zuchtbecken bzw. Verteilung erweitern.)

\subsection{Phase 5: Usage}
    Mögliche Risiken dieser Phase: 
    - Alle vorherigen Risiken. (Maßnahme: Automatisierte Überprüfung der Ergebnisse auf regelmäßiger Basis)


\subsection{Ergebnisse des Gedankenexperiment}
Die Durchführung des Gedankenexperiments hat gezeigt, dass der Einsatz sowohl eines 
Prozessmodells, als auch einer darin integrierten Risikobetrachtung 
zu mehr Sicherheit bei der Projektdurchführung sorgen kann.

